\chapter{Química supramolecular}\label{ch:b3:supramolecular}

Em 1962, Charles Pedersen, procurando outra coisa, isolou um poliéter
cíclico que dissolvia o permanganato de potássio no benzeno e o tornava
violeta. O anel tinha se enrolado em torno do íon potássio e escondido sua
carga atrás de uma superfície de hidrocarboneto, arrastando o permanganato como
contraíon. A química além da molécula, a química dos agregados mantidos
por forças mais fracas que as ligações covalentes, começou com aquela solução violeta.
Este capítulo mede essas forças, explica por que alguns \omterm{def:b3:supramolecular:supramolecular}{hospedeiros} escolhem tão bem seus \omterm{def:b3:supramolecular:supramolecular}{hóspedes},
mostra como se medem as constantes de ligação e termina com moléculas
enfiadas e entrelaçadas em máquinas.

\begin{recall}
O volume do 1.º ano tratou das forças intermoleculares, da ligação de hidrogênio, da solvatação,
das moléculas hidrofóbicas, dos complexos, das constantes de formação e dos quelatos; o volume do
2.º ano, do potencial químico, das energias de Gibbs e do ajuste por mínimos quadrados.
O \cref{ch:b3:advanced-nmr} descreveu a troca rápida e lenta e a coalescência,
o \cref{ch:b3:rate-theories}, a \omterm{thm:b3:rate-theories:eyring}{equação de Eyring}, e o
\cref{ch:b3:partition-functions}, a entropia translacional de um gás.
\end{recall}

\section{Interações não covalentes}

\begin{definition}[Química supramolecular]\label{def:b3:supramolecular:supramolecular}
A \emph{química supramolecular}\index{quimica supramolecular@química supramolecular} é a química
dos agregados de moléculas mantidos por interações não covalentes. Em um
\emph{complexo hospedeiro--hóspede}\index{complexo hospedeiro--hospede@complexo hospedeiro--hóspede}, o \emph{hospedeiro}\index{hospedeiro}
é a molécula maior, com uma cavidade ou uma fenda e sítios de ligação convergindo para
ela, e o \emph{hóspede}\index{hospede@hóspede}, a molécula ou o íon menor que ele liga.
\end{definition}

As interações são as do volume do 1.º ano, íon--dipolo, dipolo--dipolo,
ligações de hidrogênio e forças de London, mais algumas com nomes próprios.

\begin{definition}[Interações não covalentes]\label{def:b3:supramolecular:interactions}
O \emph{empilhamento $\pi$}\index{empilhamento $\pi$} é a atração entre anéis aromáticos planos
dispostos face a face, em geral deslocados ou inclinados. Uma \emph{interação
cátion--$\pi$}\index{interacao cation--pi@interação cátion--$\pi$} é a atração entre um cátion e
a face rica em elétrons de um anel aromático. Uma \emph{ligação de halogênio}\index{ligacao de halogenio@ligação de halogênio}
é a atração entre a ponta pobre em elétrons de um átomo de halogênio ligado
covalentemente, oposta à sua ligação, e uma base de Lewis. O \emph{efeito
hidrofóbico}\index{efeito hidrofobico@efeito hidrofóbico} é a tendência das superfícies apolares de se
juntarem na água, impulsionada sobretudo pela liberação das moléculas de água ordenadas
em torno delas.
\end{definition}

As interações íon--dipolo e cátion--$\pi$ são as mais fortes e dependem
muito do solvente, que compete pelo íon; as ligações de hidrogênio vêm a seguir,
depois o empilhamento $\pi$ e as \omterm{def:b3:supramolecular:interactions}{ligações de halogênio}, depois as forças de London, fracas uma a uma, mas
somadas sobre grandes superfícies de contato. Na água, o \omterm{def:b3:supramolecular:interactions}{efeito hidrofóbico} muitas vezes
domina: um \omterm{def:b3:supramolecular:supramolecular}{hospedeiro} construído para ligar um \omterm{def:b3:supramolecular:supramolecular}{hóspede} por ligações de hidrogênio no clorofórmio pode não
ligá-lo de forma alguma na água, cujas próprias ligações de hidrogênio competem.

\section{Reconhecimento molecular}

\begin{definition}[Reconhecimento molecular]\label{def:b3:supramolecular:recognition}
O \emph{reconhecimento molecular}\index{reconhecimento molecular} é a ligação
seletiva de um \omterm{def:b3:supramolecular:supramolecular}{hóspede} por um \omterm{def:b3:supramolecular:supramolecular}{hospedeiro}. Ele se apoia na \emph{complementaridade}\index{complementaridade},
uma correspondência de tamanho, forma e sítios de ligação entre \omterm{def:b3:supramolecular:supramolecular}{hospedeiro} e \omterm{def:b3:supramolecular:supramolecular}{hóspede}, e é
reforçado pela \emph{pré-organização}\index{pre-organizacao@pré-organização}: um \omterm{def:b3:supramolecular:supramolecular}{hospedeiro} cujos
sítios de ligação já são mantidos na geometria do complexo antes da ligação.
\end{definition}

A constante de associação $K_a = [\mathrm{HG}]/([\mathrm H][\mathrm G])$ é uma constante de formação no sentido do
volume do 1.º ano, e $\Delta G^\circ = -RT\ln K_a$. Um \omterm{def:b3:supramolecular:supramolecular}{hospedeiro} flexível precisa congelar sua
conformação ao se ligar, o que custa entropia; um \omterm{def:b3:supramolecular:supramolecular}{hospedeiro} rígido, pré-organizado, já
pagou esse preço em sua síntese.

\begin{definition}[Efeitos quelato e macrocíclico]\label{def:b3:supramolecular:macrocyclic}
O \emph{efeito quelato}\index{efeito quelato} é a maior estabilidade de um
complexo com um ligante polidentado do que com o número equivalente de
ligantes monodentados com os mesmos átomos doadores. O \emph{efeito
macrocíclico}\index{efeito macrociclico@efeito macrocíclico} é a estabilização adicional quando o
ligante polidentado é um anel, e não uma cadeia aberta.
\end{definition}

\begin{proposition}[O efeito quelato é sobretudo entrópico]\label{prop:b3:supramolecular:chelate-entropy}
Em $\mathrm{[M(NH_3)_6]^{2+} + 3\,en \rightleftharpoons [M(en)_3]^{2+} + 6\,NH_3}$, onde en é a etano-1,2-diamina, as mesmas seis
ligações metal--nitrogênio estão presentes dos dois lados, e a reação é impulsionada pelo
aumento do número de moléculas livres.
\end{proposition}

\begin{proof}[Argumento]
As ligações formadas e rompidas são do mesmo tipo, de modo que $\Delta_rH^\circ$ é pequeno. O número de
moléculas independentes passa de quatro para sete: três graus de liberdade translacionais
a mais de uma molécula inteira, cada um valendo dezenas de joules por kelvin
e por mol em solução (\cref{ch:b3:partition-functions}), logo $\Delta_rS^\circ > 0$ e $K > 1$.
Do ponto de vista cinético, uma vez ligada uma extremidade de um quelato, a outra é mantida perto do
metal em alta concentração efetiva e se liga imediatamente se se soltar.
\end{proof}

\begin{proposition}[Compensação entalpia--entropia]\label{prop:b3:supramolecular:compensation}
Em uma série de \omterm{def:b3:supramolecular:supramolecular}{complexos hospedeiro--hóspede} aparentados, um $\Delta H^\circ$ de ligação mais negativo
costuma vir acompanhado de um $\Delta S^\circ$ mais negativo, de modo que $\Delta G^\circ$ varia menos que cada um deles (uma
observação experimental).
\end{proposition}

Um complexo mais apertado se liga mais fortemente, mas restringe mais os movimentos do \omterm{def:b3:supramolecular:supramolecular}{hospedeiro} e
do \omterm{def:b3:supramolecular:supramolecular}{hóspede}; uma ligação de hidrogênio mais forte com o \omterm{def:b3:supramolecular:supramolecular}{hóspede} muitas vezes significa uma mais fraca com
a água que ele deslocou. A energia de Gibbs, que fixa a seletividade, é a pequena
diferença entre dois termos grandes.

\begin{definition}[Hospedeiros]\label{def:b3:supramolecular:hosts}
Um \emph{éter-coroa}\index{eter-coroa@éter-coroa} é um poliéter macrocíclico, como o
18-coroa-6, $\ce{(CH2CH2O)6}$. Um \emph{criptando}\index{criptando} é um poliéter
bicíclico com duas cabeças de ponte nitrogenadas, que envolve um cátion em três
dimensões em um \emph{criptato}\index{criptato}. Uma
\emph{ciclodextrina}\index{ciclodextrina} é um oligômero cíclico de unidades de glicose, com
a forma de um cone truncado, com uma cavidade hidrofóbica e grupos hidróxi em
suas bordas. Um \emph{calixareno}\index{calixareno} é um macrociclo em forma de taça formado por
unidades de fenol unidas por pontes metileno.
\end{definition}

\begin{omfigure}
\begin{tikzpicture}[font=\small]
  % 18-crown-6 with K+
  \foreach \k in {0,...,17} {
    \pgfmathsetmacro\ang{90 + 20*\k}
    \pgfmathsetmacro\nxt{110 + 20*\k}
    \draw[black!70] (\ang:1.55) -- (\nxt:1.55);
  }
  \foreach \k in {0,3,...,15} {
    \pgfmathsetmacro\ang{90 + 20*\k}
    \node[fill=white, inner sep=0.5pt, text=omThm] (o\k) at (\ang:1.55) {O};
    \draw[omThm!60, densely dashed] (0,0) -- (\ang:1.32);
  }
  \node[circle, fill=omDef!25, draw=omDef, inner sep=1pt] at (0,0) {$\ce{K+}$};
  \node[font=\scriptsize] at (0,-2.1) {18-coroa-6 com $\ce{K+}$};
  % [2.2.2]cryptand with K+
  \begin{scope}[xshift=5.2cm]
    \node (nt) at (0,1.7) {N};
    \node (nb) at (0,-1.7) {N};
    \foreach \x/\n in {-1.25/a, 0/b, 1.25/c} {
      \node[text=omThm, fill=white, inner sep=0.5pt] (oa\n) at (\x,0.55) {O};
      \node[text=omThm, fill=white, inner sep=0.5pt] (ob\n) at (\x,-0.55) {O};
      \draw[black!70] (nt) to[out=-90+40*\x, in=90] (oa\n) -- (ob\n) to[out=-90, in=90-40*\x] (nb);
    }
    \node[circle, fill=omDef!25, draw=omDef, inner sep=1pt] at (0.62,0) {$\ce{K+}$};
    \node[font=\scriptsize] at (0,-2.4) {[2.2.2]criptando: criptato de $\ce{K+}$};
  \end{scope}
\end{tikzpicture}

{\small À esquerda: o 18-coroa-6, um anel de doze átomos de carbono (vértices) e seis átomos de
oxigênio, com um íon potássio em seu centro ligado pelos seis oxigênios. À direita: o
[2.2.2]\omterm{def:b3:supramolecular:hosts}{criptando}, três cadeias de dois oxigênios de éter entre duas cabeças de ponte
nitrogenadas, esquemático: o cátion é mantido dentro de uma gaiola tridimensional,
pelos seis oxigênios e pelos dois nitrogênios.}
\end{omfigure}

Os éteres-coroa selecionam os cátions pelo tamanho: o 18-coroa-6 liga o $\ce{K+}$ mais fortemente que o
$\ce{Na+}$, menor, que não alcança os seis oxigênios ao mesmo tempo, ou que o $\ce{Cs+}$, maior, que
fica acima do anel. Os \omterm{def:b3:supramolecular:hosts}{criptandos}, mais pré-organizados e envolvendo completamente o íon,
ligam ainda mais fortemente e selecionam com mais nitidez. Escondido dentro
do \omterm{def:b3:supramolecular:supramolecular}{hospedeiro}, o cátion arrasta seu ânion para os solventes apolares, onde o ânion,
não solvatado, torna-se muito reativo: a base da catálise por transferência de fase.

\section{Medir a ligação}

\begin{proposition}[Troca rápida]\label{prop:b3:supramolecular:fast-exchange}
Quando um \omterm{def:b3:supramolecular:supramolecular}{hospedeiro} troca entre seus estados livre e ligado muito mais depressa que a
diferença de suas frequências de ressonância, seu sinal de RMN aparece na média
$\delta_{\mathrm{obs}} = x_{\mathrm{free}}\delta_{\mathrm{free}} + x_{\mathrm{bound}}\delta_{\mathrm{bound}}$, onde os $x$ são as frações do \omterm{def:b3:supramolecular:supramolecular}{hospedeiro} em cada estado.
\end{proposition}

\begin{proof}
Cada molécula de \omterm{def:b3:supramolecular:supramolecular}{hospedeiro} visita os dois estados muitas vezes durante a medida; sua
frequência de precessão, média no tempo, é a média ponderada das duas, com
pesos iguais às frações de tempo passadas em cada um, que no equilíbrio são as
frações de moléculas em cada estado (\cref{ch:b3:advanced-nmr}).
\end{proof}

\begin{theorem}[Isoterma 1:1 exata]\label{thm:b3:supramolecular:isotherm}
Para $\mathrm{H + G \rightleftharpoons HG}$ com concentrações totais $H_0$ e $G_0$ e constante de associação $K$,
\[
  [\mathrm{HG}] = \frac{1}{2}\Bigl[\Bigl(H_0 + G_0 + \frac1K\Bigr) - \sqrt{\Bigl(H_0 + G_0 + \frac1K\Bigr)^2 - 4H_0G_0}\,\Bigr].
\]
\end{theorem}

\begin{proof}
Com $c = [\mathrm{HG}]$, $K = c/((H_0 - c)(G_0 - c))$, isto é, $c^2 - (H_0 + G_0 + 1/K)c + H_0G_0 = 0$. Das duas raízes,
só a menor fica abaixo de $H_0$ e de $G_0$ (seu produto é $H_0G_0$ e sua soma
excede $H_0 + G_0$): é a de sinal menos.
\end{proof}

\begin{omfigure}
\begin{tikzpicture}
\begin{axis}[name=a, width=0.48\linewidth, height=5.5cm, xmin=0, xmax=6, ymin=0, ymax=1.05,
  axis lines=left, xlabel={equivalentes de hóspede, $G_0/H_0$}, ylabel={fração de hospedeiro ligada},
  tick label style={font=\scriptsize}, label style={font=\small},
  ]
\addplot[omMeth, thick] table[x=equiv, y=K4] {figdata/out/bachelor-3-binding-titration.dat};
\addplot[omDef, thick] table[x=equiv, y=K3] {figdata/out/bachelor-3-binding-titration.dat};
\addplot[omThm, thick] table[x=equiv, y=K2] {figdata/out/bachelor-3-binding-titration.dat};
\node[font=\scriptsize, omMeth, anchor=north] at (axis cs:4,0.94) {$K = 10^4$};
\node[font=\scriptsize, omDef, anchor=north] at (axis cs:4,0.68) {$K = 10^3$};
\node[font=\scriptsize, omThm, anchor=north] at (axis cs:4,0.27) {$K = 10^2$};
\end{axis}
\begin{axis}[at={(a.east)}, anchor=west, xshift=1.4cm, width=0.48\linewidth, height=5.5cm,
  xmin=0, xmax=1, ymin=0, ymax=1.05, axis lines=left,
  xlabel={fração molar de hospedeiro}, ylabel={$[\text{complexo}]/c_{\text{total}}$},
  tick label style={font=\scriptsize}, label style={font=\small}]
\addplot[omDef, thick] table[x=x, y=HG] {figdata/out/bachelor-3-binding-titration-b.dat};
\addplot[omThm, thick] table[x=x, y=HG2] {figdata/out/bachelor-3-binding-titration-b.dat};
\draw[black!35, dashed] (axis cs:0.5,0) -- (axis cs:0.5,0.55) (axis cs:0.3333,0) -- (axis cs:0.3333,0.38);
\node[font=\scriptsize, omDef, anchor=south] at (axis cs:0.5,0.52) {HG};
\node[font=\scriptsize, omThm, anchor=south] at (axis cs:0.25,0.36) {$\mathrm{HG_2}$};
\end{axis}
\end{tikzpicture}

{\small À esquerda: fração de um \omterm{def:b3:supramolecular:supramolecular}{hospedeiro} ligada durante uma titulação com $H_0 = \qty{1.0}{mmol/L}$ para três
constantes de associação (em $\unit{L.mol^{-1}}$): quanto maior $K$, mais nítida a quebra em um equivalente;
quando $KH_0 \gg 1$, a curva diz apenas que $K$ é grande. À direita: \omterm{def:b3:supramolecular:job}{gráficos de Job}, a
concentração de complexo a concentração total constante de \omterm{def:b3:supramolecular:supramolecular}{hospedeiro} mais \omterm{def:b3:supramolecular:supramolecular}{hóspede},
com máximo em uma fração de \omterm{def:b3:supramolecular:supramolecular}{hospedeiro} de 1/2 para um complexo 1:1 e de 1/3 para um 1:2 (constantes
do modelo).}
\end{omfigure}

\begin{method}[Constante de ligação a partir de uma titulação por RMN]\label{met:b3:supramolecular:nmr-titration}
\begin{enumerate}
  \item Mantenha fixa a concentração de \omterm{def:b3:supramolecular:supramolecular}{hospedeiro} $H_0$ e adicione \omterm{def:b3:supramolecular:supramolecular}{hóspede}, de 0 a vários
        equivalentes; registre um sinal do \omterm{def:b3:supramolecular:supramolecular}{hospedeiro} em cada ponto (troca rápida).
  \item Modelo: $\delta = \delta_{\mathrm{free}} + (\delta_{\mathrm{bound}} - \delta_{\mathrm{free}})[\mathrm{HG}]/H_0$, com $[\mathrm{HG}]$ dada pela isoterma exata.
  \item Ajuste $K$ e $\delta_{\mathrm{bound}}$ por mínimos quadrados não lineares, minimizando a soma dos quadrados dos
        resíduos; obtenha as incertezas-padrão a partir da curvatura dessa soma
        (a matriz de covariância).
  \item Verifique se os resíduos mostram uma tendência (estequiometria errada) e se a
        fração ligada cobre aproximadamente de 20 a 80\,\% (os dados então restringem $K$).
\end{enumerate}
\end{method}

\begin{definition}[Gráfico de Job]\label{def:b3:supramolecular:job}
Um \emph{gráfico de Job}\index{grafico de Job@gráfico de Job} (método das variações contínuas) é um gráfico da
concentração de um complexo, ou de um sinal proporcional a ela, em função da
fração molar de um componente, a concentração total constante dos
dois.
\end{definition}

\begin{proposition}[Máximo de um gráfico de Job]\label{prop:b3:supramolecular:job}
Para um único complexo $\mathrm{H_nG_m}$, o \omterm{def:b3:supramolecular:job}{gráfico de Job} tem máximo na fração molar de \omterm{def:b3:supramolecular:supramolecular}{hospedeiro}
$x = n/(n + m)$, qualquer que seja a constante de ligação.
\end{proposition}

\begin{proof}
Com concentração total $T$, $H_t = xT$, $G_t = (1 - x)T$, complexo $c$, livres $[\mathrm H] = H_t - nc$,
$[\mathrm G] = G_t - mc$, e $c = \beta[\mathrm H]^n[\mathrm G]^m$. Derivando em relação a $x$ no máximo, onde
$\dd c/\dd x = 0$: $0 = \beta\bigl(n[\mathrm H]^{n-1}[\mathrm G]^m T - m[\mathrm H]^n[\mathrm G]^{m-1}T\bigr)$, logo $n[\mathrm G] = m[\mathrm H]$. Substituindo,
$n(G_t - mc) = m(H_t - nc)$, logo $nG_t = mH_t$, isto é, $n(1 - x) = mx$: $x = n/(n + m)$.
\end{proof}

\begin{method}[Realizar um gráfico de Job]\label{met:b3:supramolecular:job}
\begin{enumerate}
  \item Prepare soluções-estoque equimolares de \omterm{def:b3:supramolecular:supramolecular}{hospedeiro} e de \omterm{def:b3:supramolecular:supramolecular}{hóspede}.
  \item Misture-as em proporções $x$ de 0 a 1 a volume total constante.
  \item Meça uma propriedade proporcional ao complexo (uma absorbância corrigida
        das espécies livres, ou $x\,\Delta\delta$ em RMN).
  \item Leia a estequiometria a partir da posição do máximo.
\end{enumerate}
\end{method}

\begin{definition}[Calorimetria de titulação isotérmica]\label{def:b3:supramolecular:itc}
A \emph{calorimetria de titulação isotérmica}\index{calorimetria de titulacao isotermica@calorimetria de titulação isotérmica}
(ITC) mede o calor liberado ou absorvido a cada injeção de uma solução de \omterm{def:b3:supramolecular:supramolecular}{hóspede}
em uma solução de \omterm{def:b3:supramolecular:supramolecular}{hospedeiro} mantida a temperatura constante.
\end{definition}

\begin{proposition}[Calor por injeção]\label{prop:b3:supramolecular:itc}
Desprezando a diluição, o calor da $i$-ésima injeção em uma célula de volume $V_0$ é
$q_i = \Delta H^\circ V_0\bigl([\mathrm{HG}]_i - [\mathrm{HG}]_{i-1}\bigr)$.
\end{proposition}

\begin{proof}
O calor é a entalpia de reação vezes a quantidade de complexo formada durante a
injeção, e essa quantidade é a variação de $[\mathrm{HG}]$ vezes o volume da célula.
\end{proof}

\begin{method}[Ler uma isoterma de ITC]\label{met:b3:supramolecular:itc}
\begin{enumerate}
  \item Trace o calor por mol de \omterm{def:b3:supramolecular:supramolecular}{hóspede} injetado em função da razão molar.
  \item A altura do patamar antes da saturação dá $\Delta H^\circ$; a razão molar no
        ponto de inflexão dá a estequiometria; a inclinação dá $K$.
  \item Ajuste os três com a isoterma exata; então $\Delta G^\circ = -RT\ln K$ e
        $T\Delta S^\circ = \Delta H^\circ - \Delta G^\circ$.
\end{enumerate}
\end{method}

Um único experimento de ITC dá, assim, toda a assinatura termodinâmica da
ligação, entalpia e entropia, que uma titulação por RMN ou por UV--visível só dá
por meio de uma série de temperaturas.

\begin{inthelab}[Uma titulação por RMN]
Uma solução do \omterm{def:b3:supramolecular:supramolecular}{hospedeiro} em um solvente deuterado, de cerca de um milimol por litro,
é colocada em um tubo de RMN e seu espectro é registrado. Uma solução concentrada do
\omterm{def:b3:supramolecular:supramolecular}{hóspede} preparada na própria solução do \omterm{def:b3:supramolecular:supramolecular}{hospedeiro}, de modo que a concentração de \omterm{def:b3:supramolecular:supramolecular}{hospedeiro} permaneça
constante, é adicionada em pequenas porções com microsseringa; após cada adição, o
tubo é agitado, equilibrado à temperatura da sonda, e o espectro é
registrado. O deslocamento de um sinal do \omterm{def:b3:supramolecular:supramolecular}{hospedeiro} que se move é lido em cada ponto e
ajustado.
\end{inthelab}

\section{Hospedeiros}

As \omterm{def:b3:supramolecular:hosts}{ciclodextrinas}, produzidas por \omterm{def:b3:complex-kinetics:enzyme}{enzimas} a partir do amido, têm seis, sete ou oito unidades de
glicose ($\alpha$, $\beta$, $\gamma$). Sua cavidade hidrofóbica acolhe as partes apolares dos \omterm{def:b3:supramolecular:supramolecular}{hóspedes} na
água: elas solubilizam fármacos, transportam fragrâncias e retiram colesterol das
membranas celulares. Os \omterm{def:b3:supramolecular:hosts}{calixarenos} e os cucurbiturilos, em forma de abóbora,
macrociclos rígidos de unidades de glicolurila com portais revestidos de carbonilas, ligam cátions e
cátions orgânicos muito fortemente na água.

\begin{omfigure}
\begin{tikzpicture}[font=\small]
  \draw[black!70, fill=omDef!8] (-1.9,1.2) -- (-1.2,-1.2) arc[start angle=180, end angle=360, x radius=1.2, y radius=0.3]
    -- (1.9,1.2);
  \draw[black!70, fill=omDef!18] (0,1.2) ellipse[x radius=1.9, y radius=0.45];
  \draw[black!40, dashed] (-1.2,-1.2) arc[start angle=180, end angle=0, x radius=1.2, y radius=0.3];
  \node[draw=omThm, thick, regular polygon, regular polygon sides=6, minimum size=8mm, inner sep=0pt] at (0,0.2) {};
  \draw[omThm, thick] (0,0.62) -- (0,2.0);
  \node[font=\scriptsize, omThm, anchor=west] at (0.1,2.0) {hóspede};
  \node[font=\scriptsize, anchor=west] at (2.0,1.2) {borda larga: OH secundários};
  \node[font=\scriptsize, anchor=west] at (1.3,-1.25) {borda estreita: OH primários};
  \node[font=\scriptsize, anchor=east] at (-1.75,0) {cavidade hidrofóbica};
\end{tikzpicture}

{\small Uma \omterm{def:b3:supramolecular:hosts}{ciclodextrina}, desenhada como o cone truncado que ela forma, com um \omterm{def:b3:supramolecular:supramolecular}{hóspede}
aromático em sua cavidade. Os grupos hidróxi das duas bordas tornam o exterior
solúvel em água; o interior, revestido de grupos C--H e de oxigênios glicosídicos, é
apolar.}
\end{omfigure}

\section{Automontagem e máquinas moleculares}

\begin{definition}[Automontagem]\label{def:b3:supramolecular:self-assembly}
A \emph{automontagem}\index{automontagem} é a organização espontânea e reversível
de componentes em uma estrutura definida, sob o controle da
informação contida nos próprios componentes: suas formas e seus sítios de
ligação.
\end{definition}

Íons metálicos com ângulos de coordenação fixos e ligantes em ponte rígidos se montam
em quadrados, gaiolas e cápsulas em uma única etapa, porque cada ligação pode se abrir
e se fechar até que se atinja a estrutura mais estável, sem tensão nem sítios livres.
O pareamento de bases no DNA e o enovelamento das proteínas são a mesma ideia
em escala maior.

\begin{definition}[Moléculas mecanicamente entrelaçadas]\label{def:b3:supramolecular:interlocked}
Uma \emph{molécula mecanicamente entrelaçada}\index{molecula mecanicamente entrelacada@molécula mecanicamente entrelaçada}
tem partes que não estão ligadas umas às outras, mas não podem ser separadas
sem romper uma ligação covalente. Um \emph{rotaxano}\index{rotaxano} é um anel
enfiado em um eixo com tampões volumosos nas duas extremidades; um
\emph{catenano}\index{catenano} é formado por anéis entrelaçados.
\end{definition}

\begin{definition}[Síntese por molde]\label{def:b3:supramolecular:template}
Uma \emph{síntese por molde}\index{sintese por molde@síntese por molde} usa uma
interação temporária, muitas vezes com um íon metálico, para manter componentes reativos na
geometria necessária a uma reação que, de outro modo, seria improvável, como o
fechamento de um anel em torno de outro anel.
\end{definition}

\begin{omfigure}
\begin{tikzpicture}[font=\small, >=Stealth]
  % step 1: two crescents around Cu+
  \draw[omDef, line width=2pt] (0.3,0) arc[start angle=0, end angle=300, radius=0.75];
  \draw[omThm, line width=2pt] (-0.3,0) arc[start angle=180, end angle=-120, radius=0.75];
  \fill[omProp] (0,0) circle (0.13);
  \node[font=\scriptsize] at (0,-1.35) {o $\ce{Cu+}$ mantém dois};
  \node[font=\scriptsize] at (0,-1.7) {ligantes abertos cruzados};
  \draw[->, thick] (1.4,0) -- (2.3,0) node[midway, above, font=\scriptsize] {fechar};
  % step 2: closed interlocked rings with Cu
  \begin{scope}[xshift=3.6cm]
    \draw[omDef, line width=2pt] (-0.35,0) circle (0.75);
    \draw[omThm, line width=2pt] (0.35,0) circle (0.75);
    \draw[white, line width=5pt] ([shift={(52:0.75)}]-0.35,0) arc[start angle=52, end angle=72, radius=0.75];
    \draw[omDef, line width=2pt] ([shift={(48:0.75)}]-0.35,0) arc[start angle=48, end angle=76, radius=0.75];
    \fill[omProp] (0,0) circle (0.13);
    \node[font=\scriptsize] at (0,-1.35) {catenato};
  \end{scope}
  \draw[->, thick] (5.0,0) -- (5.9,0) node[midway, above, font=\scriptsize] {$-\ce{Cu+}$};
  \begin{scope}[xshift=7.2cm]
    \draw[omDef, line width=2pt] (-0.35,0) circle (0.75);
    \draw[omThm, line width=2pt] (0.35,0) circle (0.75);
    \draw[white, line width=5pt] ([shift={(52:0.75)}]-0.35,0) arc[start angle=52, end angle=72, radius=0.75];
    \draw[omDef, line width=2pt] ([shift={(48:0.75)}]-0.35,0) arc[start angle=48, end angle=76, radius=0.75];
    \node[font=\scriptsize] at (0,-1.35) {[2]catenano};
  \end{scope}
  % rotaxane
  \begin{scope}[xshift=10.2cm]
    \draw[black!70, line width=1.6pt] (-1.1,0) -- (1.1,0);
    \fill[black!60] (-1.25,0) circle (0.2) (1.25,0) circle (0.2);
    \draw[omThm, line width=2pt] (0,0) ellipse[x radius=0.18, y radius=0.55];
    \node[font=\scriptsize] at (0,-1.35) {[2]rotaxano};
  \end{scope}
\end{tikzpicture}

{\small \omterm{def:b3:supramolecular:template}{Síntese por molde} de um \omterm{def:b3:supramolecular:interlocked}{catenano} (esquemática): um íon cobre(I) mantém dois
ligantes derivados da fenantrolina cruzados em ângulo reto em sua coordenação
tetraédrica; fechar cada um deles em um anel dá dois anéis entrelaçados em torno do
metal, e a remoção do cobre deixa o \omterm{def:b3:supramolecular:interlocked}{catenano} livre. À direita: um \omterm{def:b3:supramolecular:interlocked}{rotaxano}, um
anel preso em um eixo por dois tampões.}
\end{omfigure}

\begin{definition}[Máquinas moleculares]\label{def:b3:supramolecular:machine}
Uma \emph{máquina molecular}\index{maquina molecular@máquina molecular} é um agregado cujos
componentes se movem uns em relação aos outros de modo controlado em resposta a um
estímulo (luz, uma mudança redox, uma mudança de pH) e que pode ser levado a realizar trabalho.
\end{definition}

Em uma lançadeira molecular, um anel no eixo de um \omterm{def:b3:supramolecular:interlocked}{rotaxano} fica em uma de duas estações; um
estímulo que enfraquece sua ligação à primeira o envia para a segunda, e o
estímulo inverso o traz de volta. Os motores rotativos movidos a luz baseados em
alquenos sobrecarregados giram em um único sentido, alternando isomerizações
fotoquímicas e etapas térmicas que não podem ocorrer em sentido inverso. A velocidade de
deslocamento entre estações é medida por RMN a temperatura variável, a partir da
coalescência dos sinais das duas formas (\cref{ch:b3:advanced-nmr}).

\begin{safety}{\ghs[0.62cm]{GHS07}}
% ledger: ghs4:18C6
18-Coroa-6: nocivo se ingerido, irritante para a pele, os olhos e as vias respiratórias.
Os éteres-coroa transportam sais metálicos para solventes orgânicos e através da pele; são
pesados e manipulados com luvas.
\end{safety}

\begin{history}[A química supramolecular e suas máquinas]
Os éteres-coroa de Charles Pedersen (1967), os \omterm{def:b3:supramolecular:hosts}{criptandos} de Jean-Marie Lehn e os \omterm{def:b3:supramolecular:supramolecular}{hospedeiros}
pré-organizados de Donald Cram fundaram a área; os três dividiram o Prêmio Nobel de
Química de 1987. Jean-Pierre Sauvage preparou \omterm{def:b3:supramolecular:interlocked}{catenanos} com alto rendimento pelo
molde de cobre em 1983, Fraser Stoddart construiu lançadeiras de \omterm{def:b3:supramolecular:interlocked}{rotaxano}, e Ben
Feringa, um motor rotativo movido a luz em 1999; eles dividiram o Prêmio Nobel de
Química de 2016 pelo projeto e pela síntese de \omterm{def:b3:supramolecular:machine}{máquinas moleculares}.
\end{history}

\section{Exercícios}

\begin{exercise}[$\star$]\label{exo:b3:supramolecular:1}
Dê a principal interação em: $\ce{K+}$ no 18-coroa-6; dois anéis benzênicos face a face;
$\ce{K+}$ acima de um anel benzênico; iodopentafluorobenzeno com piridina; um par
guanina--citosina.
\end{exercise}

\begin{exercise}[$\star$]\label{exo:b3:supramolecular:2}
Calcule $\Delta G^\circ$ a \qty{298}{K} para um \omterm{def:b3:supramolecular:supramolecular}{complexo hospedeiro--hóspede} com $K_a = \qty{1.0e4}{L.mol^{-1}}$.
\end{exercise}

\begin{exercise}[$\star$]\label{exo:b3:supramolecular:3}
No metanol, $\log K = 6.1$ para $\ce{K+}$ e 4.4 para $\ce{Na+}$ com o 18-coroa-6 (dados do exercício).
Calcule a seletividade e a diferença de $\Delta G^\circ$.
\end{exercise}

\begin{exercise}[$\star$]\label{exo:b3:supramolecular:4}
Um \omterm{def:b3:supramolecular:job}{gráfico de Job} tem máximo em uma fração molar de \omterm{def:b3:supramolecular:supramolecular}{hospedeiro} de 0.33. Qual é a estequiometria?
\end{exercise}

\begin{exercise}[$\star\star$]\label{exo:b3:supramolecular:5}
Um \omterm{def:b3:supramolecular:supramolecular}{hospedeiro} a \qty{2.0}{mmol/L} com $K = \qty{1500}{L.mol^{-1}}$, $\delta_{\mathrm{free}} = \qty{3.560}{ppm}$ e $\delta_{\mathrm{bound}} = \qty{3.720}{ppm}$ recebe um
equivalente de \omterm{def:b3:supramolecular:supramolecular}{hóspede}. Calcule a fração ligada e o deslocamento observado.
\end{exercise}

\begin{exercise}[$\star\star$]\label{exo:b3:supramolecular:6}
Para o níquel(II), $\log\beta_6 = 8.6$ com a amônia e $\log\beta_3 = 18.3$ com a etano-1,2-diamina (dados do
exercício). Calcule a constante e $\Delta G^\circ$ a \qty{298}{K} de
$\ce{[Ni(NH3)6]^2+} + 3\,\mathrm{en} \rightleftharpoons \ce{[Ni(en)3]^2+} + \ce{6NH3}$, e diga o que a impulsiona.
\end{exercise}

\begin{exercise}[$\star\star$]\label{exo:b3:supramolecular:7}
Um experimento de ITC a \qty{298}{K} dá $K = \qty{2.0e5}{L.mol^{-1}}$ e $\Delta H^\circ = \qty{-40}{kJ/mol}$. Calcule $\Delta G^\circ$, $T\Delta S^\circ$ e
$\Delta S^\circ$.
\end{exercise}

\begin{exercise}[$\star\star$]\label{exo:b3:supramolecular:8}
Explique o papel do cobre(I), e de sua geometria tetraédrica, na síntese de \omterm{def:b3:supramolecular:interlocked}{catenanos} de
Sauvage. Como o cobre é removido no fim?
\end{exercise}

\begin{exercise}[$\star\star$]\label{exo:b3:supramolecular:9}
Os dois sinais de uma lançadeira de \omterm{def:b3:supramolecular:interlocked}{rotaxano}, separados por \qty{120}{Hz}, coalescem a \qty{300}{K} (dados do
exercício). Calcule a constante de velocidade de troca na coalescência e a \omterm{def:b3:rate-theories:activation}{energia de Gibbs
de ativação} do deslocamento.
\end{exercise}

\begin{exercise}[$\star\star\star$]\label{exo:b3:supramolecular:10}
Um fármaco tem solubilidade intrínseca de \qty{0.10}{mmol/L} e forma um complexo 1:1 com uma
\omterm{def:b3:supramolecular:hosts}{ciclodextrina}, $K = \qty{500}{L.mol^{-1}}$ (dados do exercício). Mostre que a solubilidade total em uma
solução de \omterm{def:b3:supramolecular:hosts}{ciclodextrina} total $C_t$ é $S_0 + KS_0C_t/(1 + KS_0)$, e calcule-a para $C_t = \qty{10}{mmol/L}$.
\end{exercise}

\begin{exercise}[$\star\star\star$]\label{exo:b3:supramolecular:11}
Mostre, a partir da isoterma exata, que quando $KH_0 \gg 1$ a fração ligada é igual ao
número de equivalentes até um e depois permanece igual a um. Por que uma titulação desse tipo
dá apenas um limite inferior para $K$?
\end{exercise}

\begin{exercise}[$\star\star\star$]\label{exo:b3:supramolecular:12}
Para o \omterm{def:b3:supramolecular:supramolecular}{hospedeiro} do \cref{exo:b3:supramolecular:5}, suponha a associação
\omterm{def:b3:rate-theories:diffusion}{controlada por difusão}, $k_{\mathrm{on}} = \qty{1e9}{L.mol^{-1}.s^{-1}}$. Calcule $k_{\mathrm{off}}$ e mostre que a troca é rápida na
escala de tempo da RMN a \qty{400}{MHz}.
\end{exercise}

\section{Problema: Benzeno violeta}

\begin{problem}[{Problema de fim de semana --- benzeno violeta: o potássio em um éter-coroa,
a extração do permanganato para um solvente orgânico, uma constante de ligação
a partir de uma titulação por RMN e por que uma quantidade catalítica de éter-coroa basta}]
\label{pb:b3:supramolecular:1}
% ledger: const:R, ghs:KMnO4, ghs:toluene, rion:K+VI, rion:Na+VI
Dados do problema. No metanol a \qty{298}{K}, $\log K = 6.1$ para $\ce{K+}$ e 4.4 para $\ce{Na+}$ com o
18-coroa-6. Extração: uma fase aquosa com \qty{0.10}{mol/L} de $\ce{KCl}$ e \qty{1.0}{mmol/L} de $\ce{KMnO4}$ é agitada
com igual volume de tolueno contendo o \omterm{def:b3:supramolecular:hosts}{éter-coroa} L; a única espécie
extraída é o par iônico $\ce{[KL]+}\ce{MnO4-}$, com
$K_{\mathrm{ex}} = [\ce{[KL]+MnO4-}]_{\mathrm{org}}/([\ce{K+}]_{\mathrm{aq}}[\ce{MnO4-}]_{\mathrm{aq}}[\mathrm L]_{\mathrm{org}}) = \qty{1.0e5}{L^2.mol^{-2}}$, e o \omterm{def:b3:supramolecular:hosts}{éter-coroa} permanece no tolueno. Titulação por RMN de um
\omterm{def:b3:supramolecular:hosts}{éter-coroa} (\qty{2.0}{mmol/L}) com um sal de sódio em metanol deuterado, um sinal de $\ce{CH2}$
(ppm) em função dos equivalentes de $\ce{Na+}$:

\begin{center}\small
\begin{tabular}{l*{11}{c}}
\hline
equiv. & 0 & 0.25 & 0.5 & 0.75 & 1.0 & 1.5 & 2.0 & 3.0 & 4.0 & 6.0 & 8.0\\
$\delta$ & 3.560 & 3.590 & 3.612 & 3.635 & 3.649 & 3.674 & 3.688 & 3.697 & 3.704 & 3.708 & 3.714\\
\hline
\end{tabular}
\end{center}
% table: binding-titration-c

\medskip\noindent\textbf{Parte I --- O complexo.}
\begin{enumerate}
  \item Desenhe o 18-coroa-6 e seu complexo de potássio.
  \item Calcule $\Delta G^\circ$ de ligação do $\ce{K+}$ e do $\ce{Na+}$.
  \item Calcule a seletividade $\ce{K+}/\ce{Na+}$.
  \item Com os raios de coordenação seis $r(\ce{K+}) = \qty{138}{pm}$ e $r(\ce{Na+}) = \qty{102}{pm}$, explique a seletividade.
  \item Por que o complexo é solúvel no tolueno, quando o $\ce{KCl}$ não é?
  \item Por que o \omterm{def:b3:supramolecular:hosts}{éter-coroa} se liga menos fortemente na água que no metanol?
\end{enumerate}

\medskip\noindent\textbf{Parte II --- Extração.}
\begin{enumerate}[resume]
  \item Escreva o equilíbrio de extração.
  \item Com $y$ a concentração extraída, escreva a equação de $y$ quando
        $[\mathrm L]_0 = \qty{1.0}{mmol/L}$.
  \item Resolva-a e dê a fração de permanganato extraída.
  \item Calcule a fração com $[\mathrm L]_0 = \qty{10}{mmol/L}$.
  \item Calcule-a com $[\mathrm L]_0 = \qty{0.10}{mmol/L}$.
  \item Por que o tolueno fica violeta, e por que o permanganato extraído é um oxidante
        mais forte que na água?
\end{enumerate}

\medskip\noindent\textbf{Parte III --- Uma titulação por RMN.}
\begin{enumerate}[resume]
  \item Por que um único sinal médio se desloca, em vez de dois sinais?
  \item Escreva o modelo de $\delta$ em função da concentração de \omterm{def:b3:supramolecular:supramolecular}{hóspede}.
  \item Ajuste $K$ e $\delta_{\mathrm{bound}}$ por mínimos quadrados; dê $K$ com sua incerteza-padrão.
  \item Dê a dispersão dos resíduos e comente.
  \item Em que faixa de equivalentes a fração ligada fica entre 20 e 80\,\%?
  \item Por que a mesma titulação não permitiria medir $K$ para o potássio, $\log K = 6.1$?
\end{enumerate}

\medskip\noindent\textbf{Parte IV --- Catálise por transferência de fase.}
\begin{enumerate}[resume]
  \item Um alqueno no tolueno é oxidado pelo permanganato extraído. Onde ocorre
        a reação?
  \item O que acontece com o \omterm{def:b3:supramolecular:hosts}{éter-coroa} após a reação, e por que uma quantidade catalítica
        basta?
  \item Por que o excesso de $\ce{KCl}$ na água é útil?
  \item Que catalisadores mais baratos fazem o mesmo trabalho na indústria?
  \item Por que hoje se usa o tolueno onde Pedersen usava o benzeno?
  \item Enuncie o resultado: a fração de permanganato extraída com
        $[\mathrm L]_0 = \qty{1.0}{mmol/L}$.
\end{enumerate}
\end{problem}
